% ── §8  Conclusion ────────────────────────────────────────────────────────────
\section{Conclusion}
\label{sec:conclusion}

We presented \xtal, a closed-loop, cross-layer, explainable tail-latency
controller for real-time streaming pipelines.
\xtal closes the gap between rich runtime telemetry and the large knob
space exposed by Kafka and Apache Flink, using an XGBoost contextual
bandit that learns online without offline profiling.
Its safety layer enforces correctness invariants independently of the
learned policy, and its SHAP-based explanations give operators a clear
account of every control decision.

In evaluation on NEXMark benchmarks:
\begin{itemize}
  \item \xtal reduces mean \pnn{99} latency by up to \textbf{64\%}
    during 10$\times$ burst events versus Static-Default.
  \item \xtal cuts SLA violation rate from \textbf{17.8\%} to \textbf{1.7\%}
    under bursty load and from \textbf{22.2\%} to \textbf{1.7\%}
    during TaskManager failure.
  \item \xtal achieves \textbf{2.3$\times$ faster recovery} after failure
    injection than hand-tuned static configuration.
  \item Control loop overhead is under \textbf{220\ms} per 10-second
    interval, imposing negligible cost on the pipeline.
\end{itemize}

\paragraph{Future Work.}
\xtal currently uses a discrete action space; continuous action spaces
with policy gradient methods are a natural extension.
Multi-job control (tuning across multiple concurrent Flink jobs sharing
a cluster) is an open problem.
Finally, extending explainability from per-decision to
\emph{counterfactual} explanations (``what would have happened if we
had not acted?'') would further improve operator trust.
